% ============================================================================
% Chapter 31: Formulas and Functions
% ============================================================================
\chapter{Formulas and Functions}
\label{ch:excel-formulas}

\begin{objectives}
  \item Write basic formulas using +, -, *, /
  \item Use functions: SUM, AVERAGE, MIN, MAX, COUNT
  \item Understand cell references in formulas
  \item Use the fill handle to copy formulas
\end{objectives}

\bytetip[tip]{This is where Excel truly shines! Formulas and functions turn your spreadsheet from a simple table into a powerful calculator. Let's unlock this superpower!}

\section{Writing Formulas}
\label{sec:formulas}
\index{MS Excel!formulas}

Every formula in Excel starts with an \textbf{equals sign} (\texttt{=}).

% --- Figure 31.1: Anatomy of a Formula ---
\begin{figure}[H]
\centering
\begin{tikzpicture}
  % Formula text
  \node[font=\Large\ttfamily\bfseries, text=NavyBlue] (formula) at (5,2) {=SUM(B2:B6)};
  
  % Bracket labels
  \draw[WarnRed, line width=1.5pt, ->] (1.6,2) -- (1.6,0.5);
  \node[below, font=\small\sffamily\bfseries, text=WarnRed, align=center] at (1.6,0.5) {\texttt{=}\\Always starts\\with equals};
  
  \draw[FunPurple, line width=1.5pt, ->] (3.2,2) -- (3.2,0.5);
  \node[below, font=\small\sffamily\bfseries, text=FunPurple, align=center] at (3.2,0.5) {\texttt{SUM}\\Function name\\(what to do)};
  
  \draw[NoteOrange, line width=1.5pt, ->] (4.4,2) -- (4.4,0.5);
  \node[below, font=\small\sffamily\bfseries, text=NoteOrange, align=center] at (4.4,0.5) {\texttt{(}\\Opening\\bracket};
  
  \draw[LabGreen, line width=1.5pt, ->] (6.0,2) -- (6.0,0.5);
  \node[below, font=\small\sffamily\bfseries, text=LabGreen, align=center] at (6.0,0.5) {\texttt{B2:B6}\\Cell range\\(B2 through B6)};
  
  \draw[FunSky, line width=1.5pt, ->] (7.6,2) -- (7.6,0.5);
  \node[below, font=\small\sffamily\bfseries, text=FunSky, align=center] at (7.6,0.5) {\texttt{)}\\Closing\\bracket};
\end{tikzpicture}
\caption{Anatomy of a formula --- every part explained}
\label{fig:formula-anatomy}
\end{figure}

\section{Essential Functions}
\label{sec:functions}
\index{MS Excel!functions}

% --- Figure 31.2: Functions Comparison Cards ---
\begin{figure}[H]
\centering
\begin{tabularx}{\textwidth}{@{}*{5}{>{\centering\arraybackslash}X}@{}}
  \begin{tcolorbox}[colback=ModuleBlue!8, colframe=ModuleBlue, width=2.4cm, height=3.6cm, valign=center, halign=center, arc=6pt]
    {\Large\faPlus}\vspace{4pt}\\
    \textbf{SUM}\vspace{2pt}\\
    {\scriptsize\texttt{=SUM(A1:A5)}}\vspace{2pt}\\
    {\scriptsize Adds all}\vspace{2pt}\\
    {\scriptsize Result: \textbf{50}}
  \end{tcolorbox}
  &
  \begin{tcolorbox}[colback=FunPurple!8, colframe=FunPurple, width=2.4cm, height=3.6cm, valign=center, halign=center, arc=6pt]
    {\Large\faDivide}\vspace{4pt}\\
    \textbf{AVERAGE}\vspace{2pt}\\
    {\scriptsize\texttt{=AVERAGE(..)}}\vspace{2pt}\\
    {\scriptsize Finds mean}\vspace{2pt}\\
    {\scriptsize Result: \textbf{10}}
  \end{tcolorbox}
  &
  \begin{tcolorbox}[colback=FunCoral!8, colframe=FunCoral, width=2.4cm, height=3.6cm, valign=center, halign=center, arc=6pt]
    {\Large\faArrowDown}\vspace{4pt}\\
    \textbf{MIN}\vspace{2pt}\\
    {\scriptsize\texttt{=MIN(A1:A5)}}\vspace{2pt}\\
    {\scriptsize Smallest}\vspace{2pt}\\
    {\scriptsize Result: \textbf{2}}
  \end{tcolorbox}
  &
  \begin{tcolorbox}[colback=LabGreen!8, colframe=LabGreen, width=2.4cm, height=3.6cm, valign=center, halign=center, arc=6pt]
    {\Large\faArrowUp}\vspace{4pt}\\
    \textbf{MAX}\vspace{2pt}\\
    {\scriptsize\texttt{=MAX(A1:A5)}}\vspace{2pt}\\
    {\scriptsize Largest}\vspace{2pt}\\
    {\scriptsize Result: \textbf{18}}
  \end{tcolorbox}
  &
  \begin{tcolorbox}[colback=NoteOrange!8, colframe=NoteOrange, width=2.4cm, height=3.6cm, valign=center, halign=center, arc=6pt]
    {\Large\faHashtag}\vspace{4pt}\\
    \textbf{COUNT}\vspace{2pt}\\
    {\scriptsize\texttt{=COUNT(..)}}\vspace{2pt}\\
    {\scriptsize Counts num}\vspace{2pt}\\
    {\scriptsize Result: \textbf{5}}
  \end{tcolorbox}
\end{tabularx}
\caption{Five essential Excel functions you must know}
\label{fig:functions-cards}
\end{figure}

% --- Table 31.1: Common Functions Reference ---
\begin{table}[H]
\centering
\caption{Common Excel functions reference}
\label{tab:functions}
\renewcommand{\arraystretch}{1.3}
\begin{tabularx}{\textwidth}{l X l l}
\toprule
\rowcolor{HeaderRow}
\textcolor{white}{\textbf{Function}} & \textcolor{white}{\textbf{What It Does}} & \textcolor{white}{\textbf{Example}} & \textcolor{white}{\textbf{Result}} \\
\midrule
\rowcolor{RowLight} SUM & Adds all values in a range & \texttt{=SUM(A1:A5)} & 50 \\
\rowcolor{RowDark} AVERAGE & Finds the mean (average) & \texttt{=AVERAGE(A1:A5)} & 10 \\
\rowcolor{RowLight} MIN & Finds the smallest value & \texttt{=MIN(A1:A5)} & 2 \\
\rowcolor{RowDark} MAX & Finds the largest value & \texttt{=MAX(A1:A5)} & 18 \\
\rowcolor{RowLight} COUNT & Counts cells with numbers & \texttt{=COUNT(A1:A5)} & 5 \\
\bottomrule
\end{tabularx}
\end{table}

\section{The Fill Handle}
\label{sec:fill-handle}
\index{MS Excel!fill handle}

% --- Figure 31.3: Fill Handle Demo ---
\begin{figure}[H]
\centering
\begin{tikzpicture}
  % Cells
  \foreach \y/\val/\style in {5/1/solid, 4/2/solid, 3/3/dashed, 2/4/dashed, 1/5/dashed, 0/6/dashed}{
    \draw[NavyBlue!50, line width=0.5pt, \style] (0,\y) rectangle (2,\y+1);
    \node[font=\small\sffamily] at (1,\y+0.5) {\val};
    \node[font=\tiny\sffamily, text=NavyBlue!60] at (-0.4,\y+0.5) {A\the\numexpr 6-\y};
  }
  
  % Fill handle square
  \fill[LabGreen] (1.85,4.85) rectangle (2.15,5.15);
  
  % Arrow
  \draw[LabGreen, line width=2pt, ->] (2,4.85) -- (2,0.5);
  \node[right=0.5cm of {(2,2.5)}, font=\small\sffamily\bfseries, text=LabGreen, align=center] {Drag down to\\fill the series!};
\end{tikzpicture}
\caption{The fill handle copies formulas and extends data patterns}
\label{fig:fill-handle}
\end{figure}

\bytetip[warn]{Remember: formulas always start with \texttt{=}. If you type \texttt{SUM(A1:A5)} without the equals sign, Excel treats it as text, not a formula! You'll see the text ``SUM(A1:A5)'' instead of the result.}

\bytetip[fun]{Excel can handle formulas with up to \textbf{8,192 characters} and \textbf{64 levels of nested functions}. Some finance professionals write formulas so complex they look like alien code!}

\begin{funfact}
The most commonly used function in Excel is \textbf{SUM}, followed by AVERAGE and VLOOKUP. Together, these three functions are used in over \textbf{90\%} of all Excel spreadsheets worldwide!
\end{funfact}

\begin{reviewqs}
  \item What symbol must every formula start with?
  \item Write the formula to add cells B1 through B10.
  \item What does the AVERAGE function do?
  \item What is the difference between MIN and MAX?
  \item What is the fill handle and how do you use it?
\end{reviewqs}

\begin{challenge}
Create a ``Test Scores'' spreadsheet with 5 students and 3 subjects. In column E, use \texttt{=SUM} to calculate each student's total. In row 8, use \texttt{=AVERAGE} for each subject's average. In row 9, use \texttt{=MAX} for the highest score. In row 10, use \texttt{=MIN} for the lowest.
\end{challenge}
